\subsection{An odd-discriminant two-unit family satisfying the actual GD conditions}
\label{sec:kt-odd-discriminant-two-unit}

The even discriminants in KT-EXPLICIT-TWO-UNIT-FAMILY cannot satisfy
the source's coprimality condition at $N=2$.  The following related
construction repairs that specific obstruction.  Its elementary unit
calculation is given in full; the real class-number estimate and the
$p$-rationality criterion are imported at the exact locators below.

\begin{theorem}[KT-ODD-DISCRIMINANT-TWO-UNIT-FAMILY]
Let $p\geq7$ be a prime with $p\equiv3\pmod4$.  For each
$s\in\{1,-1\}$ write
\[
 p+4s=m_su_s^2,\qquad m_s>0\text{ square-free},\quad u_s>0,
 \qquad D_s=-m_s,\qquad E_s=\mathbf Q(\sqrt{D_s}).       \tag{OD.1}
\]
Then $m_s\equiv3\pmod4$, $m_s\geq3$, and $D_s$ is an odd
negative fundamental discriminant prime to $p$.  Moreover
\[
 0<h(E_s)<p,\qquad
 \chi_{D_+}(p)=-1,\qquad \chi_{D_-}(p)=1.                \tag{OD.2}
\]
For every integer $t\geq0$ put
\[
 e=\frac{p-3}{2}+(p-1)t.                                \tag{OD.3}
\]
Then $e$ is positive and even, $(p-1)\nmid2e$, and both
\[
 B_{1,\chi_{D_s}},\qquad B_{1,\chi_{D_s}\omega^e}
                                                            \tag{OD.4}
\]
are $p$-adic units.  Equivalently for the second term,
$B_{e+1,\chi_{D_s}}/(e+1)$ is a $p$-adic unit.
The twists $\omega^e$ and $\omega^{2e}$ both have order
$(p-1)/2>2$.  The remaining third-unit condition is precisely
\[
 B_{1,\chi_{D_s}\omega^{2e}}\in\mathbf Z_p^\times
 \quad\Longleftrightarrow\quad
 \frac{B_{p-2,\chi_{D_s}}}{p-2}\in\mathbf Z_p^\times.     \tag{OD.5}
\]
No assertion of (OD.5) is part of this theorem.
\end{theorem}

\begin{proof}
Both integers $p+4s$ are positive, odd, and congruent to $3$ modulo
$4$.  Thus $u_s$ is odd and $m_s\equiv3\pmod4$, which also gives
$m_s\geq3$.  Since $p\nmid p+4s$, neither $m_s$ nor $u_s$ is
divisible by $p$.  The discriminant of $E_s$ is consequently the
odd fundamental discriminant $D_s=-m_s$.  Modulo $p$ we have
$D_s\equiv-4s/u_s^2$, so the two local symbols in (OD.2) follow
from $(-1/p)=-1$.

We first bound the imaginary class number.  For $m_s>3$, the field
$E_s$ has exactly two roots of unity.  The finite class-number
formula of Benmerieme--Movahhedi, Theorem~4.1, gives
\[
 h(E_s)=\frac{1}{2-\chi_{D_s}(2)}
       \sum_{r=1}^{(m_s-1)/2}\chi_{D_s}(r)
 \leq\frac{m_s-1}{2}\leq\frac{p+3}{2}<p.                \tag{OD.6}
\]
Here $2-\chi_{D_s}(2)$ is either $1$ or $3$, and each summand
is at most $1$.  For $m_s=3$ the same formula, with six roots of
unity, gives $h(-3)=1$.  Positivity of class numbers proves (OD.2).
The imaginary quadratic formula
$B_{1,\chi_{D_s}}=-2h(E_s)/w(E_s)$, with $w(E_s)\in\{2,6\}$,
therefore proves the first unit in (OD.4), since $p\geq7$.

For the second unit consider the real quadratic field and its unit
\[
 d_s=pm_s,\qquad F_s=\mathbf Q(\sqrt{d_s}),\qquad
 \alpha_s=\frac{p+2s+u_s\sqrt{d_s}}2
          =\frac{p+2s+\sqrt{p(p+4s)}}2.                 \tag{OD.7}
\]
The integer $d_s$ is square-free and congruent to $1$ modulo $4$,
so it is the fundamental discriminant of $F_s$.  The numerator
coefficients in (OD.7) are odd, hence $\alpha_s$ is integral.
Its norm is $((p+2s)^2-p(p+4s))/4=1$, and $\alpha_s>1$.

In fact $\alpha_s$ is the fundamental unit of $F_s$.
No integral unit can have norm $-1$: writing it as
$(x+y\sqrt{d_s})/2$ would give $x^2-d_sy^2=-4$, which is
impossible modulo $p$ because $(-1/p)=-1$.  Let $\varepsilon>1$
be the fundamental unit, of norm $1$.  Write
$\varepsilon=(x+y\sqrt{d_s})/2$, with integers $x>0$ and
$y\geq1$.  If $\alpha_s=\varepsilon^n$ for $n\geq2$, the
strict monotonicity of $z+z^{-1}$ for $z>1$ gives
\[
 p+2s=\operatorname{Tr}(\alpha_s)
 \geq\operatorname{Tr}(\varepsilon^2)
 =x^2-2=d_sy^2+2\geq3p+2>p+2,                         \tag{OD.8}
\]
a contradiction.  Thus $n=1$.

There is a unique prime $v$ of $F_s$ above $p$, with ramification
index two and residue field $\mathbf F_p$.  In the completion
$F_{s,v}$ put $\pi=\sqrt{d_s}$ and normalize $v_\pi(\pi)=1$;
then $v_\pi(p)=2$ and $u_s$ is a local unit.  Directly from
(OD.7),
\[
 s\alpha_s-1=\frac{s(p+u_s\pi)}2,
 \qquad v_\pi(s\alpha_s-1)=1.                           \tag{OD.9}
\]
Suppose $\alpha_s=z^p$ in $F_{s,v}$.  The element $z$ is a unit
and its residue is $s$, because its $p$-th power has that residue
in $\mathbf F_p$.  Thus $sz=1+w$ with $v_\pi(w)\geq1$, and
$(1+w)^p=s\alpha_s$.  For $1\leq j<p$,
$v_\pi(\binom pj w^j)\geq2+j\geq3$, while
$v_\pi(w^p)\geq p\geq7$.  The binomial expansion therefore
gives $v_\pi((1+w)^p-1)\geq3$, contradicting (OD.9).
This proves that the fundamental unit is not a local $p$-th power,
without a logarithm convention or an additional weight hypothesis.

Le's Theorem~(a) applies to every real quadratic field and gives
\[
 h(F_s)\leq\left\lfloor\frac{\sqrt{d_s}}2\right\rfloor
 \leq\frac{\sqrt{p(p+4)}}2<p.                            \tag{OD.10}
\]
Its published correction changes only the attribution of one
reference, not the theorem.  Consequently $p\nmid h(F_s)$.
Benmerieme--Movahhedi, Corollary~2.6, now proves that $F_s$ is
$p$-rational; its hypotheses require only $p\geq5$ in the present
real quadratic setting.  Their Corollary~2.2 consequently gives
\[
                    L_p(1,\chi_{F_s})\in\mathbf Z_p^\times.
                                                            \tag{OD.11}
\]

Set $n_0=(p-1)/2$.  The primitive product of the quadratic
characters $\chi_{F_s}$ and $\omega^{n_0}=\chi_{-p}$ is
$\chi_{D_s}$: indeed $(pm_s)(-p)=-m_sp^2$.
This is equality of primitive characters, equivalently of Galois
characters.  Equation~(3.2) of Benmerieme--Movahhedi therefore gives
\[
 L_p(1,\chi_{F_s})\equiv
       -\frac{B_{n_0,\chi_{D_s}}}{n_0}\pmod p.          \tag{OD.12}
\]
As an imprimitive Dirichlet function the product still omits the
$p$ Euler factor.  That factor is
$1-\chi_{D_s}(p)p^{n_0-1}\equiv1\pmod p$, since $n_0\geq3$,
so (OD.12) has exactly the displayed normalization.  The square-free
integer $-m_s$ satisfies the hypotheses preceding that equation:
it is negative, prime to $p$, and different from $-p$.

Finally apply KT-FIELD-CHARACTER-REDUCTION, since
$e+1\equiv n_0\pmod{p-1}$ and $2e+1\equiv p-2\pmod{p-1}$.
It gives the second unit in (OD.4) and the equivalence (OD.5),
including weights divisible by $p$ and weights larger than $p$.
Writing $p=4a+3$ gives $e\equiv2a\pmod{p-1}$ and
$\gcd(4a+2,2a)=2$.  Hence the two twists have the asserted common
odd order $(p-1)/2$, and $2e\equiv-2\pmod{p-1}$ proves the
nondivisibility assertion.
\end{proof}

\begin{corollary}[KT-ODD-DISCRIMINANT-GD]
For every field $E_s$, weight $e$, and prime $p$ of the theorem,
the actual source conditions $\mathrm{GD}$ hold at $N=2$ and
$\psi=1$, after choosing the usual embeddings and coefficient
data.  Either prescribed splitting behavior at $p$ is available.
\end{corollary}
\begin{proof}
We check all twelve fields of
\texttt{HeightAndVanishing/Counts.lean:9--22}.
The prime is odd, the weight is positive, and $N=2$ is even.
The character $\psi=1$ has order one, its prime-to-$p$ and
$p$-power parts are trivial, and its parity agrees with even $e$.
The field is imaginary quadratic.  By (OD.1)--(OD.2),
$p\nmid |D_s|h(E_s)$ and $\gcd(|D_s|,2)=1$.
The final genericity condition is $(p-1)\nmid2e$, already proved.

For the remaining field \texttt{eta}, the source requires that
the intersection of the kernels of the residual cyclotomic
character and the trivial Dirichlet character is not contained
in $G_{E_s}$.  Choose any prime $q\mid m_s$; such a prime exists
because $m_s\geq3$, and $q\ne p$.  The quadratic character
$\chi_{D_s}$ is nontrivial on inertia at $q$, since $q$ divides
the fundamental discriminant.  The residual cyclotomic character
at $p$ is unramified there, and the Dirichlet character is trivial.
An inertia element on which $\chi_{D_s}$ is $-1$ therefore
belongs to the specified intersection but not to $G_{E_s}$.
This is precisely the remaining noncontainment condition.
\end{proof}

\paragraph{Arithmetic scope.}
For these actual GD fields, the already proved
KT-ORDINARY-RANK2-UNIT-EXCLUSION applies with $a=e$ to every
representation satisfying that theorem's explicit residual,
ordinary-quotient, and away-$p$ conditions.  Thus the constructed
ordinary lower-weight $\mathrm{GL}_2$ source branch is excluded
where those conditions hold.  The actual Type~II exclusion uses
$a=2e$ and still requires the third unit (OD.5).
Neither the third unit nor an exclusion of all genuinely rank-three
sources follows from this two-unit theorem.  No additional local
conditions away from $p$ are prescribed.

\paragraph{Exact primary locators and check scope.}
Maohua Le, \emph{Upper bounds for class numbers of real quadratic
fields}, Acta Arithmetica 68 (1994), 141--144,
DOI \texttt{10.4064/aa-68-2-141-144}: Theorem~(a) on printed
p.~141; its complete proof is in Section~3, pp.~142--143.
The complete four-page paper was read.  The numerical initial
range and the analytic lemmas cited there are imported, not
independently recomputed.  Le's one-page correction, Acta
Arithmetica 72 (1995), p.~399,
DOI \texttt{10.4064/aa-72-4-399}, was read completely; it corrects
the author of reference~[2] to B.~Oriat and makes no mathematical
change.  Publisher PDFs are saved with retrieval metadata under
\texttt{evidence/main\_gate/literature/odd\_discriminant/}.

Y.~Benmerieme and A.~Movahhedi, \emph{Multi-quadratic $p$-rational
number fields}, Journal of Pure and Applied Algebra 225 (2021),
106657, DOI \texttt{10.1016/j.jpaa.2020.106657}; version
arXiv:2007.04864v1, 9 July 2020,
\url{https://arxiv.org/pdf/2007.04864v1}:
Corollary~2.2 and preceding derivation, PDF pp.~4--5;
Proposition~2.3, Proposition~2.4, and Corollary~2.6 with their
intervening derivation, pp.~5--6; equation~(3.2) and its full
context, p.~12; Theorem~4.1, p.~14.
The authors mention the field $\mathbf Q(\sqrt{p(p-4)})$ on
p.~14, without supplying the proof used here.  The apparently
missing display after the last ``same method'' sentence on p.~15
is also blank in the rendered PDF; no unstated theorem is inferred
from it.  The present proof supplies both signs directly.
